\section{Evaluation}
\label{sec:eval}

\noindent\textbf{Experimental Setup}
\begin{figure}
\begin{tabular}{ | m{5em} | m{17em} | } 
  \hline
  Frontend & 2-way superscalar, 256 entries BTB, 8 entries RAS, tournament predictor \\
  \hline
  Reorder Buffer &  64 entries, 2-way insert/commit \\
  \hline
  Reservation Stations & 2 stations for ALU (16 entries), 1 for FPU (16 entry), 1 for memory (16 entries) \\
  \hline
  L1 (I/D) &  32KB, 8 ways, 8 outstanding requests \\
  \hline
  Load-Store Unit & 24 LdQ entries, 14 StQ entries, 4 SB entries\\
  \hline
  LLC (L2) & 1MB, 10 cycles latency, 16 ways, max 16 outstanding requests \\
  \hline
  Memory &  120 cycles latency (24 outstanding requests) \\
  \hline
\end{tabular}
    \caption{Citadel's microarchitectural configuration.}
    \label{tab:mi6param}
\end{figure}
Citadel's processor is a modified out-of-order 2-core based on Riscy-OOO~\cite{zhang2018composable} and MI6~\cite{bourgeat2019mi6}.
The parameters used for the cores and the memory subsystem are given in Figure~\ref{tab:mi6param}, and kept small to keep the synthesis and simulation time small.
Citadel runs on an AWS EC2 F1 instance, equipped with a Virtex UltraScale+ FPGA. 
On this FPGA board, Citadel can be clocked at $30$ MHz, which is similar to the baseline design.
We use the RISC-V tool chain with GCC 12.2.0.
We compile all software with the -O3 flag. 
We run Linux Kernel 6.2.0.
All performance numbers are collected by running the benchmarks on the FPGA using performance counters and an average of ten runs. 

\noindent\textbf{Implementation Data and TCB Breakdown}
\begin{table}
    \centering
    \begin{tabular}{|p{2.35cm}|r|r|}
        \hline
         Component & LOC & Size\\
         \hline\hline
         Processor \cite{bourgeat2019mi6} (+ modification) & $\sim60$ KLOC (+ 250 ) & N/A \\
         \hline
         Bootloader & 5236 & 55KB\\
         \hline
         Security Monitor & 9287 & 97KB\\
         \hline
         \rowcolor{LightRedbis}
         Mini-SM & 2301 & 5.6KB\\
         \hline
         \rowcolor{LightRed}
         Linux & $\sim20$ MLOC & 157MB\\
         \hline
         \rowcolor{LightRed}
         SM Kernel Mod. & 507 & 505KB\\
         \hline
         \hline
         TCB Software & 14523 & 152KB \\
         \hline
    \end{tabular}
    \vspace{0.40cm}
    \caption{Software Component Size and TCB Breakdown. Mini-SM is a subset of the SM code. 
    Lines in red correspond to elements excluded from the TCB.}
    \label{tab:tcbloc}
\end{table}
Number of lines of code (LOC) and binary size of Citadel's components are presented in Table \ref{tab:tcbloc}.
Due to their limited functions, the SM and the Bootloader can be kept small and the total software Trusted Code Base (TCB) is less than 15K LOC.
This helps to minimize potential bugs and reduce the software attack surface.
It can also help to make future formal verification efforts tractable.
Final TCB size needs to include the enclave's binary, which depends on the application. 

\subsection{Evaluating Secure Shared Memory}
\label{sec:evalsharedmemory}

\noindent\textbf{Microbenchmark on \texttt{memcpy}}
We evaluate two different setups, one with the function as is (only using SpecSafe) and the other one when minimally annotated for Burst mode (full C code in supplemental material and assembly snippet in Figure \ref{fig:memcpyasm}).
We normalize the results to using \texttt{memcpy} to and from private memory (not using our mechanisms) in Table \ref{tab:resultmemcpy}. 
We observe that out-of-the-box, SpecSafe shows dramatic overhead, especially when the data is not cached in the LLC with a 2231\% overhead.
On the other end, we show that Burst mode is quite efficient and reduces the overhead to 245\% and 62\% depending on whether the data is cached in the LLC.
We also try different configurations where we start with a cold LLC (all loads will go to main memory), start with a warm LLC (the OS has accessed it previously or the data just arrived from the network for instance) or an \emph{insecure} setup, where we have enabled the L1 for shared memory accesses.
Our results confirm that access latency has a huge impact on the performance of SpecSafe, which is mitigated by Burst Mode.
Not shown in the table, we observe that when moving data from private memory to shared memory, the overhead is 48\% in all cases (both mechanisms and with cold/warm LLC).
Indeed, store operations to shared memory are already pipelined with SpecSafe and Burst mode does not improve the performance.
We also measure the overhead of straight-line speculation to be around 33\% by enabling Burst mode while using \texttt{memcpy} to and from private memory.
Finally, looking at performance counters, we observe that the ~20\% overhead left can be attributed to address translations, which are already cached when performing private memory copy.

\noindent\textbf{L1-D Bypass}
These results also highlight the reasonable performance impact of bypassing the L1-D for shared memory accesses, especially when using Burst mode with a warmed up LLC.
Note that we can also make a qualitative argument for bypassing the L1-D as shared memory accesses usually do not involve re-access (computing on shared memory could be insecure as the access pattern is leaked) and getting the output data to the LLC as fast as possible makes a lot of sense as we want it to be accessed by another security domain.
\begin{table}[h]
    \centering
    \begin{tabular}{|c|c|c|}
    \cline{2-3}
    \multicolumn{1}{c|}{} & \textbf{SpecSafe} & \textbf{Burst Mode} \\ \hline
    \textbf{No LLC warm-up} & 23.317 & 3.458 \\ \hline
    \textbf{LLC warm-up} & 3.665 & 1.621 \\ \hline
    \textcolor{red}{\textbf{LLC warm-up + L1}} & 2.466 & 1.556 \\ \hline
    \end{tabular}
    \caption{\texttt{memcpy} micro-benchmark. Copy of 1500 bytes from shared memory to private memory. Normalized over \texttt{memcpy} from an to private memory with speculation enabled. We experiment with warming up the LLC and enabling the L1 for shared memory accesses (insecure).}
    \label{tab:resultmemcpy}
    \newcolumntype{C}[1]{>{\centering\arraybackslash}m{#1}}
     \begin{NiceTabular}{|C{2.8cm}|C{1.8cm}|C{1.8cm}|}[hvlines, corners=NW]
        \CodeBefore
        %\rectanglecolor[gray]{0.9}{1-1}{1-3}
        %\rectanglecolor[gray]{0.9}{1-5}{2-5}
        \rowcolor{mygreen!35}{2}
        %\rowcolor{mygreen!35}{3}
        %\rowcolor{red!35}{4}
        \Body
            \multicolumn{1}{c|}{} & \textbf{Signature} & \textbf{Hash} \\
            
            \Block{}{\textbf{Citadel\\(SpecSafe + Burst)}} & 1.027 & 1.464 \\
            \Block{}{\textbf{SpecSafe Only}} & 1.194 & 4.973 \\
            \Block{}{\textbf{\textcolor{red}{Burst Only}}} & 1.654 & 4.959 \\
    \end{NiceTabular}
    \vspace{0.40cm}
    \caption{Normalized runtime for signing and hashing three thousand random messages sent by the OS to an enclave. Runtime is normalized to the insecure baseline where the library is directly linked to the testbench program. We compare Citadel with two configurations where we are not using burst mode or insecurely using burst mode everywhere.
    }
    \label{tab:resultcrypto}
\end{table}

\noindent\textbf{Cryptographic Library}
We now evaluate our mechanisms for secure shared memory on a realistic application by wrapping a cryptographic library in an enclave (see details in Section \ref{sec:evalcrypto}).
We measure the time it takes for the enclave to perform $3000$ cryptographic operations (signature or cryptographic hash) over random inputs provided by the OS.
We compare three setups in Table \ref{tab:resultcrypto}.
First, our final design that combines spec-safe shared memory + burst mode enabled around performance sensitive shared memory burst sequences (Citadel).
Then Citadel without annotating code for shared memory burst sequences (Spec-Safe Shared Memory only).
Finally, another (insecure) configuration where we annotate all the code to be run in burst mode (Burst only).
We report the runtime normalized to the insecure baseline (directly linked library and no LLC partitioning).
We observe an overhead of $2.7\%$ on the signature benchmark and $46.4\%$ for the hashing benchmark.
This is because the hashing application accesses shared memory much more frequently compared to the signing application that is more computationally bounded.
This difference between the benchmarks is also highlighted by our other results, where not annotating the code for burst shared accesses has a $397.3\%$ overhead for the hash application and a $19.4\%$ overhead for the signing application.
These results show that the burst mode improves performance for code segments with burst shared memory accesses.

\noindent\textbf{Burst Mode in the General Case}
On our benchmarks, we observe that the performance overhead goes up to $65.4\%$ for the signing benchmark and $395.9\%$ for the hashing benchmark when (insecurely) activating burst mode for the whole enclave (Table \ref{tab:resultcrypto}). 
This overhead is due to the TLB being trashed by straight-line misspeculations across function contexts.
Instead of speculating a jump to the return address using the Return Address Stack (RAS) at the end of the function, the processor will speculatively run the prologue of the next function in memory.
We observe that many of these functions will try to deference pointers passed as arguments, triggering translations for arbitrary values in this erroneous context.
This will evict useful cached translations from the TLB causing an excessive slowdown later on.
This effect was already documented (in a nonsecurity context) in the early 2000s \cite{onur1,onur2,onur3}.
We confirm this hypothesis by instrumenting the binary and observing that when the RAS is reactivated (but not other predictors), the number of TLB misses, page walks and the overall execution time goes down significantly.
Note that we do not address this performance issue, but instead use burst mode for performance-critical code snippets, which also makes security analysis tractable (see Section \ref{sec:secannalysis}).

\subsection{Evaluating the End-to-End Platform}

\noindent\textbf{Securing a Cryptographic Library}
\label{sec:evalcrypto}
We write a wrapper around an ED25519 RISC-V library \cite{ed25519github}.
We add a queue to shared memory, so untrusted applications can send requests to the library to generate keys (that will stay in enclave memory), or to sign messages using previously generated keys.
In total, we need to implement 334 LOC of C, which is to be compared to the 4,514 LOC of the crypto library (7\% overhead).
With the enclave's code added, the total TCB is now 19,231 LOC.
The final crypto-enclave binary is 71KB.

We compare performance between three scenarios.
In each of these scenarios, we measure the number of cycles required to sign randomly generated messages.
We split \texttt{1MB} of randomly generated data into $\approx 3000$ packets of various lengths following the arbitrary Internet Mix distribution \cite{imix}.
In the first setup, which corresponds to our insecure baseline, we do not use any enclave as we directly link the cryptographic library to our user-level application.
In the second scenario, called ``multi-core'', the enclave constantly runs on a dedicated core and programs running on other cores can forward requests through the shared memory mechanism.
In the third scenario, termed ``re-entry'', the enclave is re-entered every single time an application calls the cryptographic library on the same core as that of the application.
We see very little difference (2.9\%) between the first two setups, and a more important 27.6\% slowdown between the first and third scenario.
Indeed, our cryptographic algorithm is compute-heavy, making the communication overhead negligible in the second scenario.
Nevertheless, context-switching in and out of an enclave is expensive as it requires flushing the microarchitectural state, which makes the slowdown more important in the re-entry scenario.
These two secure solutions highlight an interesting tradeoff between resource occupation (one core and/or a DRAM region) and performance.

\noindent\textbf{Private Inference}
\label{sec:evalprivateinference}
We evaluate the following scenario where a remote user wants to perform a private inference task on a neural network evaluated on an untrusted server.
We place a simple neural network model inside of one of our enclaves.
We use onnx2c \cite{onnx2cgithub} to convert a ONNX file describing a handwritten digit recognition model (trained on MNIST) to standard C code and run private inference inside our enclave.
The user uses our remote attestation mechanism to verify the enclave identity and establish a shared secret with it.
The user is then able to send 1024 encrypted requests to the enclave that are passed along to the enclave by the untrusted OS.
Upon reception, the enclave decrypts the ciphertext and performs the inference task on the private input.
It then sends back an encrypted output to the remote user.
In total, we need to implement 227 lines of straightforward C, which is to be compared to the 1536 LOC of the neural network and the 410 LOC of the AES library and the 4386 LOC of ED25519 we also include
for encryption and key exchange.
With the enclave code added, the TCB is now 21K LOC and the final enclave binary measures 54KB.

We now compare the performance between our setup and an insecure baseline where the same benchmark (including decryption / encryption) is directly evaluated outside an enclave and where the LLC is not partitioned.
In both setups, runtime is split 55\% and 45\% between AES and inference.
We see negligible (0.76\%) overhead for running a model inside the enclave, with all of it coming from the necessity to copy the ciphertext in enclave private memory before performing decryption.
This is not surprising as the rest of the computation happens in enclave private memory and should not suffer from any slowdown.
This doesn't account for the time required to set up the enclave, which is a one-time cost but represents 23.3\% of the cost of performing 1024 inferences.

\noindent\textbf{MicroPython Runtime}
\label{sec:evalmicropython}
We do a minimal port of MicroPython \cite{microphython} inside an enclave.
The goal is to provide an attested Python runtime, with a signed transcript of the interactions.
It serves as a small benchmark to demonstrate the expressivity and flexibility of the platform.
We start from the ARM port and only need to modify a few of the MicroPython runtime functions.
For instance, we redirect the console to print to shared memory and need to write a bit of assembly for the \texttt{setjmp} and \texttt{longjmp} functions used for exception handling.
Moreover, a hash of the standard output is extended on each character print and the hash is signed on the enclave exit using a cryptographic library and the enclave's own private keys (obtained through the attestation mechanism).
In total, we modify 128 LOC of MicroPython for the port to our platform, and add the 5K LOC cryptographic library. 
With the enclave's code added, the TCB is now 65,174 LOC, most of it (70\%) being MicroPython.
We test our attested runtime on a simple Python program: a simple primality test.
We successfully run the test, and verify the signature of the transcript.

\subsection{Evaluating LLC Dynamic Partitioning}
\label{sec:evalllcpartitioning}

We evaluate the advantage of dynamically partitioning the LLC on our cryptographic enclave benchmark.
We compare three different LLC configurations:
an insecure baseline where the LLC is not partitioned, a secure configuration where the LLC is statically partitioned into uniform partitions of 16KB (MI6 scheme) and a last configuration where we use our scheme to partition the LLC non-uniformly.
For the non-uniform configuration, we split the 1MB of LLC as follows: we give 256KB to the enclave, 16KB to the SM, 128KB to the OS, 128KB to the region used for shared memory, 128KB to the region storing page tables, and 1KB each to every other region.
We observe a 67.8\% overhead for the uniform partitioning and 0.08\% overhead for the non-uniform partitioning, highlighting the benefits of dynamic partitioning.
This is especially important as untrusted applications also benefit from it.
We also measure the time it takes to modify LLC partitioning.
We measure that the operation takes on average 1M cycles.
This is a very costly operation, but it only happens very infrequently.
Among these, 98\% of the cycles are spent flushing the LLC.
As a result, changing the layout of fewer regions is much faster: we measure an average of 2.5K cycles to flush a single LLC set containing 16 ways. 
The bigger the number of LLC set indexes affected, the longer the operation will take.

\subsection{Takeaways for Commercial Systems}
Our performance evaluation gives us insights on the potential impact of hardening commercial systems against side channels and transient execution attacks.
We demonstrated that hardware with both bare-metal strong isolation primitives and a controlled speculation mode, akin to our Burst mode, enables system developers to create libraries and abstractions to interact with shared memory with good performance, without taking the usual security risks related to speculated executions.
These mechanisms might be even more relevant when considering commercial-grade cache-hierarchy with higher latencies.
We envision that the same kind of controlled speculation mode is implementable in commercial machines, to give most of the benefits of speculation for easy-to-verify mostly-straight-line code when enclaves access shared memory, with a simple and user-understandable speculation model. Note that enclave memory accesses can have unrestricted speculation.

While implementing such controlled speculation mode and hardware primitives on modern industrial processors would be more involved than for our prototype, we emphasize the small relative complexity of the features compared to the rest of the design (250 LOC out of 60K LOC).
We also showed that the typical cost of LLC partitioning for isolation can be greatly mitigated with a simple design that allows dynamic resizing of the partitions.
We believe that all performance optimizations presented will also improve commercial systems, albeit with different quantitative benefits.